\chapter{Contexte général du projet}

\section*{Introduction}
Ce chapitre situe le stage dans son environnement. Il présente d'abord l'organisme
d'accueil et le projet AMR-X, puis décrit les briques logicielles existantes au début du
stage et leurs limites. À partir de cet existant, la problématique et le cahier des charges
sont définis. Le chapitre se termine par l'environnement de travail et la démarche adoptée.

\section{Présentation de l'organisme d'accueil}
\acompleter{Présenter ici l'organisme d'accueil : activité, date de création,
implantation, équipe dans laquelle le stage a été réalisé.}

\section{Le projet AMR-X}
AMR-X (\emph{Autonomous Modular Robot -- Extended}) est un projet de plateforme mobile
modulaire : une base robotisée commune reçoit des modules fonctionnels interchangeables
selon l'application visée (logistique d'entrepôt, milieu hospitalier, etc.). Le dépôt du
projet regroupe la définition produit, les études de marché, la mécanique,
l'électronique, la partie robotique ROS~2 et les outils web.

La partie robotique cible l'environnement Ubuntu~24.04, ROS~2 Jazzy et Gazebo Harmonic.
Elle comprend notamment la description du robot, sa simulation dans plusieurs mondes
(entrepôt, hôpital), la cartographie par SLAM et la navigation autonome avec Nav2. Le
projet est mené par plusieurs équipes de stagiaires qui travaillent sur des parties
différentes du dépôt ; le travail présenté ici devait donc s'intégrer à l'existant sans le
reconstruire.

\section{Existant au début du stage}

\subsection{Simulation et navigation}
Le robot AMR-X est simulé dans Gazebo Harmonic. Le monde retenu pour ce stage est
l'entrepôt \code{warehouse\_harmonic}, dans lequel le robot apparaît à une position de
départ vérifiée, $(-1{,}2639 ; -3{,}9049)$~m. Le robot est équipé de deux LiDAR dont les
mesures sont fusionnées sur un topic unique, d'une caméra et d'une centrale inertielle.

La navigation repose sur Nav2 : localisation AMCL sur une carte d'occupation enregistrée
au préalable, planificateur global NavFn et contrôleur local MPPI. Une destination peut
être envoyée à Nav2 depuis RViz ou par l'action ROS~2 \code{/navigate\_to\_pose}.

\subsection{Dashboard web}
Le dashboard est une application web de supervision composée de deux parties :
\begin{itemize}
  \item un \textbf{frontend} React (Vite) organisé en pages : vue d'ensemble, robots,
    missions, carte en direct (\emph{Live map}), alertes, modules, téléopération,
    paramètres et gestion des utilisateurs ;
  \item un \textbf{backend} FastAPI avec SQLAlchemy et Alembic, qui gère les comptes (avec
    approbation par un administrateur), les robots, les missions et les alertes dans une
    base de données relationnelle.
\end{itemize}
La page \emph{Live map} affiche le monde SDF choisi sous forme de plan 2D (SVG) ou de vue
isométrique 3D. Les mondes sont extraits automatiquement des fichiers SDF du dépôt dans un
manifeste JSON (six mondes). Une connexion rosbridge existait déjà pour certaines pages.

\subsection{Limites de l'existant}
L'analyse de l'existant a fait apparaître les limites suivantes :
\begin{itemize}
  \item la \emph{Live map} n'affichait pas la position ROS du robot : les robots visibles
    provenaient de données de démonstration ou de la base de données ;
  \item en cas d'erreur de l'API, la page basculait implicitement sur les robots de
    démonstration, sans que l'utilisateur puisse distinguer des données réelles ;
  \item la vue 3D appliquait un décalage arbitraire de $(-12 ; -8)$ et animait un robot et
    un trajet décoratifs sans lien avec la simulation ;
  \item les boutons de la page \emph{Missions} enregistrent des missions dans la base mais
    ne pilotent pas Nav2 ;
  \item aucun chemin ne permettait d'envoyer une destination au robot depuis le navigateur
    ni d'en connaître le résultat.
\end{itemize}

\begin{figure}[htbp]
\centering
\begin{tikzpicture}[node distance=1.2cm and 1.6cm]
  \node[web,minimum width=4cm] (dash) {Dashboard web\\\scriptsize React + FastAPI};
  \node[web,below=0.5cm of dash,minimum width=4cm,fill=gray!10] (db)
    {Base de données\\\scriptsize robots, missions};
  \node[ros,right=3.6cm of dash,minimum width=3.6cm] (nav) {Nav2\\\scriptsize AMCL, NavFn, MPPI};
  \node[sim,below=0.5cm of nav,minimum width=3.6cm] (gz) {Gazebo Harmonic\\\scriptsize robot AMR-X simulé};
  \draw[fl] (dash) -- (db);
  \draw[fl] (nav) -- (gz);
  \draw[fl] (gz.east) -- ++(0.5,0) |- node[lbl,pos=0.25]{capteurs} (nav.east);
  \draw[thick,dashed,red!70!black] (dash.east) -- node[lbl,text=red!70!black]{aucun lien\\
    position / commande} (nav.west);
  \node[font=\scriptsize\itshape,below=0.2cm of db] {affichage : démo ou base de données};
  \node[font=\scriptsize\itshape,below=0.2cm of gz] {commande : RViz uniquement};
\end{tikzpicture}
\caption{Briques existantes au début du stage.}
\label{fig:existant}
\end{figure}

\section{Problématique}
Le dashboard et la simulation formaient deux mondes séparés. La problématique du stage
peut se formuler ainsi : \emph{comment permettre à un opérateur de commander la navigation
du robot AMR-X simulé depuis le dashboard web existant, tout en lui garantissant que la
position et l'état affichés correspondent réellement à ceux du robot ?}

Plusieurs difficultés en découlent. Le navigateur communique en WebSocket et ne parle pas
directement ROS~2. Une commande envoyée à travers le réseau peut être perdue, dupliquée ou
arriver alors que le robot exécute déjà une autre destination. Le plan affiché dans le
navigateur est construit à partir du monde SDF de Gazebo, alors que Nav2 raisonne dans le
repère de sa carte. Enfin, l'interface doit rester honnête lorsque les données deviennent
indisponibles, par exemple si la simulation est en pause ou si la connexion est perdue.

\section{Cahier des charges}

\subsection{Besoins fonctionnels}
Le tableau~\ref{tab:besoins} résume les besoins fonctionnels retenus.

\begin{table}[htbp]
\centering\small
\begin{tabularx}{\linewidth}{@{}lX@{}}
\toprule
\textbf{Réf.} & \textbf{Besoin} \\
\midrule
BF1 & Afficher en direct la position et l'orientation du robot simulé sur la \emph{Live map} (vues 2D et 3D). \\
BF2 & Choisir une destination par clic sur la carte 2D ou par saisie de X, Y (m) et du cap (°). \\
BF3 & Envoyer la destination à Nav2 et obtenir un acquittement explicite de la demande. \\
BF4 & Suivre l'état de la navigation : envoi, navigation, distance restante, résultat final. \\
BF5 & Annuler la navigation en cours et afficher la confirmation d'annulation. \\
BF6 & Refuser les destinations invalides (hors carte, occupées, non finies, mauvais monde) avec un message explicite. \\
\bottomrule
\end{tabularx}
\caption{Besoins fonctionnels.}
\label{tab:besoins}
\end{table}

\subsection{Besoins non fonctionnels}
\begin{itemize}
  \item \textbf{Fidélité} : ne jamais afficher de position ou de résultat inventé ; le
    succès provient uniquement du résultat Nav2.
  \item \textbf{Sûreté des commandes} : un seul but actif à la fois, pas de renvoi
    automatique d'une commande dont le résultat est incertain.
  \item \textbf{Détection des données périmées} : bloquer les commandes si l'état n'a pas
    été reçu depuis 2,5~s.
  \item \textbf{Intégration} : réutiliser le dashboard, Nav2 et la simulation existants
    sans les reconstruire ; conserver l'authentification existante.
  \item \textbf{Reproductibilité} : pouvoir lancer et vérifier l'ensemble de la chaîne
    par des commandes documentées.
\end{itemize}

\subsection{Périmètre et limites}
Le périmètre a été volontairement restreint à un robot, un serveur de navigation et un
monde actif à la fois. Le travail porte sur la simulation : la passerelle refuse les
destinations lorsque le temps simulé n'est pas utilisé, afin d'éviter toute commande
d'un robot physique. La pause et la reprise d'une navigation, l'arbitrage avec d'autres
sources de commande (RViz, téléopération) et l'authentification des commandes au niveau de
rosbridge sont hors périmètre. Les navigations lancées depuis la \emph{Live map} ne sont
pas enregistrées dans la base de données.

\section{Environnement de travail}
Le tableau~\ref{tab:env} présente les outils utilisés.

\begin{table}[htbp]
\centering\small
\begin{tabular}{@{}ll@{}}
\toprule
\textbf{Élément} & \textbf{Version / outil} \\
\midrule
Système d'exploitation & Ubuntu 24.04 \\
Intergiciel robotique & ROS~2 Jazzy \\
Simulateur & Gazebo Harmonic \\
Navigation & Nav2 (AMCL, NavFn, MPPI) \\
Pont web & rosbridge\_server 2.7.1 (WebSocket, port 9090) \\
Frontend & React 19, Vite 8, roslib 2.1 \\
Backend & FastAPI, SQLAlchemy, Alembic \\
Environnement JavaScript & Node.js 20 \\
Gestion de version & Git, GitHub \\
\bottomrule
\end{tabular}
\caption{Environnement logiciel du stage.}
\label{tab:env}
\end{table}

\section{Démarche adoptée}
Le travail a été conduit en quatre étapes :
\begin{enumerate}
  \item analyse de l'existant (dashboard, simulation, configuration Nav2) et définition du
    parcours utilisateur ;
  \item conception du protocole entre le navigateur et Nav2 et de la gestion des repères ;
  \item développement de la passerelle ROS~2 et des composants du dashboard, avec des tests
    unitaires sans ROS ;
  \item intégration sous Ubuntu, correction des problèmes observés et validation par des
    essais mesurés dans Gazebo.
\end{enumerate}

\section*{Conclusion}
Ce chapitre a présenté le projet AMR-X et l'existant au début du stage : une simulation,
une pile Nav2 et un dashboard fonctionnels mais non reliés. Le cahier des charges fixe
l'objectif de commander le robot depuis le dashboard avec des informations fidèles. Le
chapitre suivant décrit la conception retenue pour y parvenir.
